\documentclass[11pt]{article}
\usepackage{mathblatt}
% vorspann.tex – Kompetenzblatt, zusammenbau v0.7 (2026-09-28).
% Ergänzt mathblatt.sty für Kompetenzblätter, ohne sie zu ändern
% (layout-befunde.md 1–34 vom 28.09.). Wird nach \usepackage{mathblatt}
% eingelesen.
\makeatletter
\RequirePackage{newunicodechar}
% Zeichen, die Latin Modern nicht hat (Befund 20): Text aus einer
% Ersatzschrift, Mathe als Befehl; ohne Ersatzschrift auch Text als Mathe.
\newif\ifkbersatz
\IfFontExistsTF{FreeSerif}{\newfontfamily\kbersatzschrift{FreeSerif}\kbersatztrue}{%
 \IfFontExistsTF{DejaVu Serif}{\newfontfamily\kbersatzschrift{DejaVu Serif}\kbersatztrue}{%
  \IfFontExistsTF{Cambria Math}{\newfontfamily\kbersatzschrift{Cambria Math}\kbersatztrue}{%
   \IfFontExistsTF{Segoe UI Symbol}{\newfontfamily\kbersatzschrift{Segoe UI Symbol}\kbersatztrue}{}}}}
\newcommand{\kbzeichen}[2]{\ifmmode #2\else\ifkbersatz{\kbersatzschrift #1}\else\ensuremath{#2}\fi\fi}
\newcommand{\kbzeichenhoch}[1]{\ifkbersatz\ifmmode\text{\kbersatzschrift #1}\else{\kbersatzschrift #1}\fi\else ?\fi}
% Kopf- und Fußzeile (Befund 1, 2): Kopf Thema · Blattart, Fuß Kennung und Seite
\newcommand{\kbkopfzeile}[2]{\pagestyle{fancy}\fancyhf{}\renewcommand{\headrulewidth}{0pt}%
  \fancyhead[L]{\footnotesize\color{mbgrau}#1}%
  \fancyfoot[L]{\footnotesize\color{mbgrau}#2}%
  \fancyfoot[R]{\footnotesize\color{mbgrau}Seite \thepage}}
% Flattersatz überall (Befund 25)
\raggedright
% Titel des Blatts: Ich-kann-Satz, darunter Zweigzeile
\newcommand{\kbtitel}[2]{\par\noindent{\Large\bfseries\raggedright #1\par}%
  \ifblank{#2}{}{\par\smallskip\noindent{\small #2}\par}\medskip}
% Abschnitt: „Das kennst du schon“, „Zum Merken“
\newcommand{\kbabschnitt}[1]{\par\addvspace{12pt}\Needspace*{5\baselineskip}%
  \noindent{\bfseries #1}\par\nobreak\vspace{1pt}\noindent{\color{mbgrau}\rule{\linewidth}{0.4pt}}\par\nobreak}
% Hauptnummer: Titel hängt am ersten Block; jeder Block (Teilaufgabe) bricht
% nicht, passt er nicht mehr, rückt er ganz auf die nächste Seite (Befund 21).
\newif\ifkbtitel
\newsavebox{\kbbox}
\newlength{\kbhoehe}
\newlength{\kbnummerabstand}\setlength{\kbnummerabstand}{10pt}
\newlength{\kbteilabstand}\setlength{\kbteilabstand}{6pt}
\newenvironment{kbaufgabe}[1]{\stepcounter{aufgabe}\setcounter{teil}{0}%
  \gdef\kbtiteltext{\noindent\hangindent1.6em\hangafter1\makebox[1.6em][l]{\textbf{\theaufgabe.}}#1\par}%
  \global\kbtiteltrue\par\addvspace{\kbnummerabstand}}{\par}
\newenvironment{kbblock}{\par\addvspace{\kbteilabstand}\begin{lrbox}{\kbbox}%
  \begin{minipage}[t]{\linewidth}\raggedright\parskip\z@
  \ifkbtitel\kbtiteltext\vspace{2pt}\global\kbtitelfalse\fi
  \list{}{\leftmargin1.6em\labelwidth1.4em\labelsep0.2em\topsep\z@\partopsep\z@
    \parsep\z@\itemsep\z@\rightmargin\z@}\raggedright}%
  {\endlist\end{minipage}\end{lrbox}%
  \setlength{\kbhoehe}{\dimexpr\ht\kbbox+\dp\kbbox\relax}%
  \Needspace*{\kbhoehe}\noindent\usebox{\kbbox}\par}
% Paarsatz (Platz nutzen, Befund 27): zwei Hauptnummern oder zwei
% Teilaufgaben mit Zeichenaufgabe nebeneinander, als ein Block
\newcommand{\kbnummer}[1]{\stepcounter{aufgabe}\setcounter{teil}{0}%
  \noindent\hangindent1.6em\hangafter1\makebox[1.6em][l]{\textbf{\theaufgabe.}}#1\par\vspace{2pt}}
\newcommand{\kbhalb}[1]{\list{}{\leftmargin1.6em\labelwidth1.4em\labelsep0.2em\topsep\z@
  \partopsep\z@\parsep\z@\itemsep\z@}\raggedright #1\endlist}
\newcommand{\kbpaar}[2]{\par\addvspace{\kbnummerabstand}\begin{lrbox}{\kbbox}%
  \begin{minipage}[t]{0.485\linewidth}\raggedright\parskip\z@ #1\end{minipage}\hfill
  \begin{minipage}[t]{0.485\linewidth}\raggedright\parskip\z@ #2\end{minipage}\end{lrbox}%
  \setlength{\kbhoehe}{\dimexpr\ht\kbbox+\dp\kbbox\relax}\Needspace*{\kbhoehe}\noindent\usebox{\kbbox}\par}
\newcommand{\kbteilpaar}[2]{\item[]\noindent
  \begin{minipage}[t]{0.485\linewidth}\kbhalb{#1}\end{minipage}\hfill
  \begin{minipage}[t]{0.485\linewidth}\kbhalb{#2}\end{minipage}\par}
% Anweisung der Hauptnummer, eigene Zeile über der Teilaufgabe (Befund 5)
\newcommand{\kbanweisung}[1]{\item[]#1\par\vspace{2pt}}
% Kopfzeile einer Teilaufgabe: Buchstabe, Auftakt (halbfett oder Kontext),
% Prüfkennung klein rechts (Befund 3, 12, Beschluss c)
\newlength{\kbkennbreite}\setlength{\kbkennbreite}{1.9cm}
\newcommand{\kbkennung}[1]{{\footnotesize #1}}
\newcommand{\kbteil}[3]{\item[#1]\ifblank{#3}%
  {\parbox[t]{\linewidth}{\raggedright #2}}%
  {\parbox[t]{\dimexpr\linewidth-\kbkennbreite\relax}{\raggedright #2}\hfill
   \parbox[t]{\kbkennbreite}{\raggedleft\kbkennung{#3}}}\par}
% Frage in eigener Zeile (Befund 4), ein Satz je Zeile (Befund 10)
\newcommand{\kbfrage}[1]{\par\vspace{2pt}#1\par}
% Antwortfeld in eigener Zeile darunter, linksbündig (Befund 4, 8)
\newcommand{\kbantwort}[1]{\par\vspace{4pt}\noindent#1\par}
% Zwei Spalten: links Grafik/Tabelle/Aufgabe, rechts Antwort oder Rechenraum
% (Befund 27); oben bündig. Beginnt einen eigenen Listenpunkt ohne Marke.
\newcommand{\kbzweispaltig}[3][0.48]{\item[]\noindent
  \begin{minipage}[t]{#1\linewidth}\raggedright #2\end{minipage}\hfill
  \begin{minipage}[t]{\dimexpr0.98\linewidth-#1\linewidth\relax}\raggedright #3\end{minipage}\par}
% Linke Spalte mit Teilaufgabe (Buchstabe hängt links heraus wie in der Liste)
\newcommand{\kbspalte}[1]{\list{}{\leftmargin\z@\labelwidth1.4em\labelsep0.2em\topsep\z@
  \partopsep\z@\parsep\z@\itemsep\z@}\raggedright #1\endlist}
% Antwortfelder: 2,2 cm, in Punkten 1,2 cm; ohne Umbruch davor
\newcommand{\kbfeld}[1][]{\mbox{\underline{\hspace{2.2cm}}\ifblank{#1}{}{\,#1}}}
\newcommand{\kbfeldk}[1][]{\mbox{\underline{\hspace{1.2cm}}\ifblank{#1}{}{\,#1}}}
% Grafik oder Tabelle unter dem Text, linksbündig (Befund 8, 28)
\newcommand{\kbdarunter}[1]{\par\vspace{4pt}\noindent#1\par}
% Ankreuzen: je Option eine Zeile, linksbündig (Befund 6, 7)
\newcommand{\kbkreuz}[1]{\par\vspace{2pt}\noindent$\square$\ #1\par}
% Bearbeitungsraum (Befund 9): Schreibzeilen wie mathblatt (9 mm)
\newcommand{\kbraum}[1]{\schreibzeilen{#1}}
% Wertetabelle mit Köpfen im Textmodus (Befund 26): wie \wertetabelle der
% Vorlage, aber die Köpfe setzt der Aufruf selbst ($x$, „Zeit in h“).
\NewDocumentCommand{\kbwertetabelle}{o m m m}{%
  \def\mbkopf{}\def\mbwerte{}\def\mbleer{}\setcounter{mbn}{0}%
  \foreach \v in {#4}{\xappto\mbkopf{& $\v$}\xappto\mbleer{& }\stepcounter{mbn}\xdef\mbnn{\arabic{mbn}}}%
  \IfNoValueTF{#1}{}{\foreach \v in {#1}{\ifdefempty{\v}{\xappto\mbwerte{& }}{\xappto\mbwerte{& $\v$}}}}%
  \begingroup\renewcommand{\arraystretch}{1.3}%
  \edef\mbpre{|>{\noexpand\columncolor{mbkasten}}l|*{\mbnn}{>{\noexpand\centering\noexpand\arraybackslash}p{\noexpand\mbzell}|}}%
  \expandafter\mbtabstart\expandafter{\mbpre}\hline
  #2 \mbkopf \\ \hline
  \IfNoValueTF{#1}{\rule{0pt}{\mbzeile}#3 \mbleer \\ \hline}{#3 \mbwerte \\ \hline}
  \end{tabular}\endgroup}
% Merkkasten am Ende (Aufbau des Kompetenzblatts)
\newcommand{\kbmerk}[2]{\par\addvspace{14pt}\begin{lrbox}{\kbbox}\begin{minipage}[t]{\linewidth}%
  \noindent{\bfseries #1}\par\vspace{1pt}\noindent{\color{mbgrau}\rule{\linewidth}{0.4pt}}\par\vspace{4pt}%
  \uebersichtskasten{\raggedright #2}\end{minipage}\end{lrbox}%
  \setlength{\kbhoehe}{\dimexpr\ht\kbbox+\dp\kbbox\relax}\Needspace*{\kbhoehe}\noindent\usebox{\kbbox}\par}
% Lösungsblatt: Nummer und Buchstabe halbfett, Lösung daneben
\newcommand{\kbloesung}[2]{\par\addvspace{3pt}\noindent\hangindent2.8em\hangafter1\makebox[2.8em][l]{\textbf{#1}}#2\par}
% Zeichen dieses Blatts, die Latin Modern nicht hat
\newunicodechar{₁}{\kbzeichenhoch{₁}}
\newunicodechar{₂}{\kbzeichenhoch{₂}}
\newunicodechar{≈}{\kbzeichen{≈}{\approx}}
\makeatother

\begin{document}
\kbkopfzeile{Quadratische Gleichungen · Kompetenzblatt}{QGL-K1}
% Kompetenzblatt QGL-K1: quadratische-gleichungen e3 „p-q-Formel“ (zusammenbau v0.7, Niveau FOR)
\kbtitel{Ich kann quadratische Gleichungen mit der p-q-Formel lösen.}{ab Kl. 9, je nach Buch bis Kl. 10 · P10 ×7}
\setcounter{aufgabe}{0}

\kbabschnitt{Das kennst du schon}

% zone: quadratische-gleichungen-zone-f8-v1
\begin{kbaufgabe}{Ich kann einen Term ordnen: x² zuerst, dann x, dann die Zahl.}
\begin{kbblock}
\kbanweisung{Ordne den Term.}
\kbteil{}{{\bfseries\boldmath $3x + x^2 + 2$}}{}
\kbantwort{\kbfeld}
\end{kbblock}
\end{kbaufgabe}

% zone: quadratische-gleichungen-zone-f7-v1
\begin{kbaufgabe}{Ich kann Klammern ausmultiplizieren.}
\begin{kbblock}
\kbanweisung{Multipliziere aus.}
\kbteil{}{{\bfseries\boldmath $x \cdot (x + 4)$}}{}
\kbantwort{\kbfeld}
\end{kbblock}
\end{kbaufgabe}

% zone: quadratische-gleichungen-zone-f2-v1
\begin{kbaufgabe}{Ich kann eine Quadratwurzel ziehen.}
\begin{kbblock}
\kbanweisung{Berechne die Wurzel.}
\kbteil{}{{\bfseries\boldmath $\sqrt{49}$}}{}
\kbantwort{\kbfeld}
\end{kbblock}
\end{kbaufgabe}

\kbabschnitt{Schritt für Schritt}

% leiter: quadratische-gleichungen-e3-k3-s0-v1
\begin{kbaufgabe}{Ich kann in einer Gleichung p und q ablesen.}
\begin{kbblock}
\kbanweisung{Lies p und q mit Vorzeichen ab.}
\kbteil{}{{\bfseries\boldmath $x^2 - x - 12 = 0$}}{}
\kbantwort{$p$ = \kbfeld, $q$ = \kbfeld}
\end{kbblock}
\end{kbaufgabe}

% leiter: quadratische-gleichungen-e3-k3-s1-v1
\begin{kbaufgabe}{Ich kann eine Gleichung mit der p-q-Formel lösen.}
\begin{kbblock}
\kbzweispaltig[0.46]{\kbspalte{\kbanweisung{Löse die Gleichung mit der p-q-Formel.}\kbteil{}{{\bfseries\boldmath $x^2 + 6x + 8 = 0$}}{}\kbantwort{$x_1$ = \kbfeld, $x_2$ = \kbfeld}}}{\kbraum{2}}
\end{kbblock}
\end{kbaufgabe}

% leiter: quadratische-gleichungen-e3-k3-s2-v1
\begin{kbaufgabe}{Ich kann die p-q-Formel auch anwenden, wenn p oder q negativ ist.}
\begin{kbblock}
\kbzweispaltig[0.46]{\kbspalte{\kbanweisung{Löse die Gleichung mit der p-q-Formel.}\kbteil{}{{\bfseries\boldmath $x^2 - 4x - 12 = 0$}}{}\kbantwort{$x_1$ = \kbfeld, $x_2$ = \kbfeld}}}{\kbraum{2}}
\end{kbblock}
\end{kbaufgabe}

% leiter: quadratische-gleichungen-e3-k3-s3-v1
\begin{kbaufgabe}{Ich kann eine Gleichung erst ordnen und dann lösen.}
\begin{kbblock}
\kbzweispaltig[0.46]{\kbspalte{\kbanweisung{Löse die Gleichung.}\kbteil{}{{\bfseries\boldmath $x^2 + 5 = 6x$}}{}\kbantwort{$x_1$ = \kbfeld, $x_2$ = \kbfeld}}}{\kbraum{2}}
\end{kbblock}
\end{kbaufgabe}

% leiter: quadratische-gleichungen-e3-k3-s4-v1
\begin{kbaufgabe}{Ich kann eine Gleichung lösen, bei der vor x² nicht 1 steht.}
\begin{kbblock}
\kbzweispaltig[0.46]{\kbspalte{\kbanweisung{Löse die Gleichung.}\kbteil{}{{\bfseries\boldmath $5x^2 - 20x + 15 = 0$}}{}\kbantwort{$x_1$ = \kbfeld, $x_2$ = \kbfeld}}}{\kbraum{2}}
\end{kbblock}
\end{kbaufgabe}

% leiter: quadratische-gleichungen-e3-k3-s5-v1
\begin{kbaufgabe}{Ich kann mit der p-q-Formel rechnen, wenn p ungerade ist.}
\begin{kbblock}
\kbzweispaltig[0.46]{\kbspalte{\kbanweisung{Löse die Gleichung.}\kbteil{}{{\bfseries\boldmath $x^2 - 3x + 2 = 0$}}{}\kbantwort{$x_1$ = \kbfeld, $x_2$ = \kbfeld}}}{\kbraum{2}}
\end{kbblock}
\end{kbaufgabe}

% leiter: quadratische-gleichungen-e3-k3-s6-v1
\begin{kbaufgabe}{Ich kann erkennen, wann eine Gleichung nur eine Lösung hat.}
\begin{kbblock}
\kbzweispaltig[0.46]{\kbspalte{\kbanweisung{Löse die Gleichung.}\kbteil{}{{\bfseries\boldmath $x^2 - 10x + 25 = 0$}}{}\kbantwort{$x_1$ = \kbfeld, $x_2$ = \kbfeld}}}{\kbraum{2}}
\end{kbblock}
\end{kbaufgabe}

% leiter: quadratische-gleichungen-e3-k3-s7-v1
\begin{kbaufgabe}{Ich kann erkennen, wann eine Gleichung keine Lösung hat.}
\begin{kbblock}
\kbzweispaltig[0.46]{\kbspalte{\kbanweisung{Löse die Gleichung.}\kbteil{}{{\bfseries\boldmath $x^2 + 2x + 5 = 0$}}{}\kbantwort{Lösung: \kbfeld}}}{\kbraum{2}}
\end{kbblock}
\end{kbaufgabe}

% leiter: quadratische-gleichungen-e3-k3-s8-v1
\begin{kbaufgabe}{Ich kann an der Diskriminante sehen, wie viele Lösungen es gibt.}
\begin{kbblock}
\kbanweisung{Berechne die Diskriminante D und gib die Zahl der Lösungen an.}
\kbteil{}{{\bfseries\boldmath $x^2 - 8x + 7 = 0$}}{}
\kbfrage{Diskriminante $D = \left(\frac{p}{2}\right)^2 - q$ und Zahl der Lösungen?}
\kbantwort{$D$ = \kbfeld}\kbantwort{Lösungen: \kbfeld}
\end{kbblock}
\end{kbaufgabe}

% leiter: quadratische-gleichungen-e3-k3-s9-v1
\begin{kbaufgabe}{Ich kann Lösungen mit Wurzel angeben und runden.}
\begin{kbblock}
\kbzweispaltig[0.46]{\kbspalte{\kbanweisung{Löse die Gleichung. Runde auf zwei Stellen.}\kbteil{}{{\bfseries\boldmath $x^2 - 6x + 4 = 0$}}{}\kbantwort{$x_1$ ≈ \kbfeld, $x_2$ ≈ \kbfeld}}}{\kbraum{2}}
\end{kbblock}
\end{kbaufgabe}

% leiter: quadratische-gleichungen-e3-k3-s10-v1
\begin{kbaufgabe}{Ich kann eine Lösung durch Einsetzen prüfen.}
\begin{kbblock}
\kbzweispaltig[0.46]{\kbspalte{\kbanweisung{Setze ein und kreuze an.}\kbteil{}{{\bfseries\boldmath Probe}}{}\kbfrage{Ist $x = -2$ Lösung von $x^2 - 3x - 10 = 0$?}\kbantwort{\janein}}}{\kbraum{2}}
\end{kbblock}
\end{kbaufgabe}

% leiter: quadratische-gleichungen-e3-k3-s11-v1
\begin{kbaufgabe}{Ich kann erst Klammern auflösen, dann ordnen und lösen.}
\begin{kbblock}
\kbzweispaltig[0.46]{\kbspalte{\kbanweisung{Löse die Gleichung.}\kbteil{}{{\bfseries\boldmath $(x + 2)^2 = 5x + 10$}}{}\kbantwort{$x_1$ = \kbfeld, $x_2$ = \kbfeld}}}{\kbraum{2}}
\end{kbblock}
\end{kbaufgabe}

% leiter: quadratische-gleichungen-e3-k3-s12-v1
\begin{kbaufgabe}{Ich kann den kürzesten Lösungsweg wählen.}
\begin{kbblock}
\kbanweisung{Nenne den Weg und löse.}
\kbteil{}{{\bfseries\boldmath $x^2 - 12x = 0$}}{}
\kbfrage{Welcher Weg, welche Lösungen?}
\kbantwort{Weg: \kbfeld}\kbantwort{$x_1$ = \kbfeld, $x_2$ = \kbfeld}
\end{kbblock}
\end{kbaufgabe}

% pruefung: quadratische-gleichungen-e3-k3-s13-v3, quadratische-gleichungen-e3-k3-s13-v11, quadratische-gleichungen-e3-k3-s13-v5, quadratische-gleichungen-e3-k3-s13-v9
\begin{kbaufgabe}{Ich kann das auch in Aufgaben aus der Prüfung.}
\begin{kbblock}
\kbteil{a)}{{\bfseries\boldmath Die Funktion $g$ mit $g(x) = x^2 - 10x + 22$}}{P10 ’25}
\kbfrage{Schnittpunkte mit der x-Achse (zwei Stellen)?}
\kbzweispaltig[0.46]{\kbantwort{( \kbfeldk | 0), ( \kbfeldk | 0)}}{\kbraum{2}}
\end{kbblock}
\begin{kbblock}
\kbteil{b)}{{\bfseries\boldmath Parabel $y = x^2 - 2$ und Gerade $y = 3x + 2$}}{P10 ’24}
\kbfrage{Schnittpunkte?}
\kbzweispaltig[0.46]{\kbantwort{( \kbfeldk | \kbfeldk ), ( \kbfeldk | \kbfeldk )}}{\kbraum{3}}
\end{kbblock}
\begin{kbblock}
\kbteil{c)}{{\bfseries\boldmath Parabel $y = -x^2 + 2x + 6$}}{P10 ’23}
\kbfrage{An welchen Stellen $x$ ist $y = -2$?}
\kbzweispaltig[0.46]{\kbantwort{$x_1$ = \kbfeld, $x_2$ = \kbfeld}}{\kbraum{3}}
\end{kbblock}
\begin{kbblock}
\kbteil{d)}{{\bfseries\boldmath Parabel $y = (x - 3)^2 - 2$ und Gerade $y = -2x + 19$}}{P10 ’22}
\kbfrage{Schnittpunkte?}
\kbzweispaltig[0.46]{\kbantwort{( \kbfeldk | \kbfeldk ), ( \kbfeldk | \kbfeldk )}}{\kbraum{3}}
\end{kbblock}
\end{kbaufgabe}

\kbmerk{Zum Merken}{Normalform: x² + px + q = 0 – alles auf eine Seite, geordnet nach x², x, Zahl; die Vorzahl von x² ist eins, sonst durch sie teilen (normieren). \\ 3x² + 24x − 60 = 0 | : 3 → x² + 8x − 20 = 0 \quad p = 8, q = −20 \\ p-q-Formel: x₁,₂ = −p/2 ± √((p/2)² − q) \\ x² + 8x − 20 = 0 → x₁,₂ = −4 ± √(16 + 20) = −4 ± 6 → x₁ = 2, x₂ = −10 \\ Unter der Wurzel entscheidet: positiv → zwei Lösungen, null → eine Lösung, negativ → keine Lösung. \\ x² + 8x + 16 = 0 → −4 ± √(16 − 16) → x = −4 \quad x² + 8x + 20 = 0 → √(16 − 20): keine Lösung}
\end{document}
